\begin{textsample}{POS Dim 1 – mg – Score 28.00 – mg/mg\_em\_97.txt}  \label{ex:f1_pos_247}
4.10.2 Estratégias para o desenvolvimento do componente curricular Para o desenvolvimento das aulas de Projeto de Vida, espera -se \textbf{que} o professor utilize estratégias diversificadas para estimular os estudantes ao diálogo, levando -os à reflexão sobre si mesmos, à compreensão para com os outros e ao entendimento com relação ao mundo circundante.
É importante \textbf{que}, durante o processo de convívio e reflexão proposto por esse componente curricular, os jovens tomem conhecimento, cada vez mais, de seus direitos, deveres e, principalmente, consciência daquilo \textbf{que} os torna seres únicos.
Assim, para além de \textbf{uma} preocupação com as demandas do mundo do trabalho e de outras imposições típicas da nossa sociedade \textbf{contemporânea}, o projeto de vida do estudante deve ser composto a partir de elementos \textbf{que} respeitem suas expectativas, conhecimentos prévios, valores e possibilidades.
E, como \textbf{salienta} MACHADO ( 2006, p. 5 ), todo projeto é constituído por \textbf{uma} “ referência ao futuro, a abertura para o novo e o caráter indelegável da ação projetada. ” Os três pontos mencionados acima \textbf{revelam} o tom \textbf{que} orienta o desenvolvimento desse componente curricular, assim como alguns cuidados \textbf{que} devemos ter ao trabalhá -lo na escola.
O projeto de vida do estudante precisa ganhar contornos e diretrizes no presente, correspondendo às demandas e expectativas individuais, no sentido de \textbf{que} o jovem tem o tempo \textbf{que} ainda está por vir para a concretização de suas propostas.
Além disso, sua estruturação deve ser flexível, \textbf{visto} as variáveis ( pessoais, sociais, econômicas, políticas, dentre outras ) \textbf{que} irão ocorrer entre a elaboração do projeto e a sua efetivação – não podemos nos esquecer, por exemplo, \textbf{que} os estudantes irão experimentar também no ambiente da escola encontros e desencontros, conquistas e frustrações, gerir expectativas e decepções, dialogar, \textbf{argumentar} e, por vezes, aprender a ceder sua opinião, além de outras situações \textbf{que} irão impactar profundamente suas perspectivas e levar a revisões e reestruturações de seus planos e postura individual.
Por fim, é indiscutível o fato de \textbf{que} o projeto de vida é exclusivamente do estudante e, por isso mesmo, intransferível : deve ser pensado, arquitetado e ulteriormente executado unicamente por ele – mesmo \textbf{que} seja constituído a partir de \textbf{influências} diversas, como conselhos de familiares ou reflexões e discussões orientadas na escola.
Os professores devem estar cientes da proposta do Projeto de Vida, subsidiando os jovens em suas reflexões e planejamentos.
Para tanto, o professor deverá reconhecer as juventudes, mantendo \textbf{uma} escuta ativa e atuando a partir de \textbf{uma} comunicação não-violenta, no intuito de compreender as dúvidas, problemas, desejos e outras características \textbf{que} o estudante carrega consigo e \textbf{que} são fundamentais para elaborar \textbf{um} projeto de vida factível.
Metodologias ativas e ações pragmáticas, \textbf{focadas} em resolução de problemas e tomadas de decisão, devem dar a tônica das atividades desenvolvidas nesse componente curricular. \textbf{Ademais}, experiências junto às pessoas \textbf{que} sejam referência para os estudantes, vivências fora da escola \textbf{que} possam ressignificar suas histórias, contato com pessoas de seu bairro e cidade, acompanhamento pedagógico \textbf{que} valorize a criatividade, a consciência crítica e a capacidade de se adaptar e intervir na realidade são elementos \textbf{que} também devem estar no \textbf{horizonte} do planejamento do professor responsável pelo componente curricular Projeto de Vida.
Isso nos mostra porque o Projeto de Vida não pode ser tratado como \textbf{um} componente curricular típico da formação geral básica, tal como aqueles \textbf{que} compõem as quatro áreas do conhecimento do presente currículo.
Se a \textbf{mera} \textbf{exposição} de conteúdos já não se mostra eficaz para a formação integral do estudante na atualidade, em se tratando de Projeto de Vida essa máxima tem \textbf{um} peso ainda maior.
Como afirma VILLAS \& NONATO ( 2014 ), > o percurso da vida é marcado por começos, interrupções e recomeços, fazendo > com \textbf{que} cada \textbf{sujeito} trace trajetórias singulares.
Cada \textbf{sujeito} responde e age > diante dessas situações de maneira diferente, construindo percursos e > identidades diferentes. ( VILLAS \& NONATO, 2014, p. 15 ) \textbf{Visto} \textbf{que} o Projeto de Vida lida com liberdades distintas, resoluções em desenvolvimento, tempos e espaços em constante movimento, vale reforçar \textbf{que} o professor deve manter \textbf{uma} postura sempre ativa, imparcial e flexível diante dessas particularidades para \textbf{conseguir} apoiar de modo satisfatório os projetos de vida desenvolvidos pelos estudantes.
Professores \textbf{que} tenham \textbf{uma} sensibilidade em identificar e trabalhar dificuldades intelectuais e emocionais, exercer sua autoridade de modo amistoso, administrar conflitos e transformar forças excedentes em ações assertivas – esse é o perfil do educador desejado para ministrar o componente curricular Projeto de Vida.
É importante \textbf{lembrarmos} \textbf{que}, tal como afirma SAVATER ( 2005 ), > o professor não apenas, \textbf{talvez} nem principalmente, ensina com seus \textbf{meros} > conhecimentos científicos, mas com a arte de persuasão de sua ascendência > sobre os \textbf{que} o atendem : deve ser capaz de seduzir sem hipnotizar.
Quantas > vezes a vocação do aluno é \textbf{despertada} mais por adesão a \textbf{um} professor > preferido do \textbf{que} à matéria \textbf{que} ele transmite! > > ( p. 110 ) O professor no contexto de ensino é o \textbf{sujeito} capaz de propor o desenvolvimento das atividades afins do Projeto de Vida de acordo com a necessidade da realidade escolar e as características e demandas dos estudantes e da comunidade.
Entretanto, \textbf{frisa} -se \textbf{que} o trabalho docente se oriente por algumas temáticas importantes de \textbf{que} os jovens do \textbf{século} XXI precisam se apropriar, a fim de elaborar \textbf{um} projeto de vida \textbf{que} os preparem tanto para o momento presente quanto para os eventos futuros, contribuindo para a formação de cidadãos ativos, críticos e propositivos. \textbf{Logo}, espera -se \textbf{que} as discussões \textbf{que} \textbf{fomentam} o projeto de vida dos estudantes favoreçam \textbf{uma} associação com o conhecimento adquirido durante todo o processo de aprendizagem.
Para \textbf{instigar} as discussões sobre o Projeto de Vida é \textbf{preciso} \textbf{que} os conteúdos curriculares trabalhem de forma interdisciplinar, refletindo questões \textbf{que} vão para além do mundo do trabalho, conectando os temas estudados em sala de aula com o cotidiano do estudante. \textbf{Visto} o \textbf{que} foi exposto nos parágrafos anteriores, evidencia -se a necessidade de mobilizar esforços, atenções e comprometimentos – dentro e fora da escola – para \textbf{que} ele seja significativo.
Estudantes, famílias, professores e demais educadores das instituições de ensino constituem elementos indispensáveis para o êxito desse novo componente curricular.
Assim sendo, \textbf{um} trabalho em regime de colaboração é fundamental, visando não apenas as condições de possibilidade para \textbf{que} todos os projetos almejados pelos estudantes ganhem contornos reais, mas também resguardar a integridade de todos os envolvidos nessa proposta.
De modo a atender essa necessidade de cuidado para com o outro, a escola deve estreitar vínculos com a comunidade e consolidar parcerias com o sistema de saúde pública e de assistência social para \textbf{amparar} e proteger adequadamente estudantes e professores envolvidos com o Projeto de Vida.
Isso é \textbf{imprescindível} já \textbf{que}, diante de eventuais temas sensíveis ( violências doméstica, sexual e física, negligência, abandono etc. ) \textbf{que} podem vir à tona durante esse momento de formação das juventudes, não podemos desconsiderar a promoção da qualidade de vida, assim como a manutenção da saúde física e mental do corpo docente e \textbf{discente} da escola.
De \textbf{um} modo geral, o Projeto de Vida se destaca no Currículo Referência do Ensino Médio de Minas Gerais como espaço e tempo \textbf{oportunos} para ressignificar e valorizar as visões de mundo não apenas dos estudantes, mas também de todos os educadores \textbf{que} fortalecem as estruturas \textbf{que} sustentam as escolas mineiras.

% matched lemmas: ademais, amparar, argumentar, conseguir, contemporâneo, despertar, discente, exposição, focar, fomentar, frisar, horizonte, imprescindível, influência, instigar, lembrar, logo, mero, oportuno, preciso, que, revelar, salientar, sujeito, século, talvez, um, ver
\end{textsample}
